% !TEX root = main.tex
\section{Introduction and Collatz Sequence}\label{sec:introduction-and-collatz-sequence} % (fold)
Our story starts from the $3n + 1$ Conjecture, hereinafter referred to as the Collatz conjecture, which is defined as follows:
\begin{conjecture}[Collatz Conjecture]~\label{conjecture:original-conjecture}
Given a number $a_0 = a \in \mathbb{N}$, we have the modulor arithmetic notation
\[a_{k+1} = \left\{\begin{aligned}
    &a_{k}/2, &a_{k}\equiv 0\pmod{2}, \\
    &3a_{k}+1, &a_{k}\equiv 1\pmod{2},
\end{aligned}\right.\]
where $k \in \mathbb{N}$.
The conjecture is, for all $a_0 \in \mathbb{N}$, there's $k \ge 0$, that $a_{k} = 1$.
\end{conjecture}

First, for all $k$, there is $m \ge 0$ such that $3a_{k}+1 \equiv 0 \pmod{2^{m}}$.
We may wander the maximum value of $m$ which depends on $a_k$, what the $a_k$ depends on $a$ and $k$.
Let's define the symbol $\val(a;k+1) = \val_{2}(3a_{k}+1)$ as count of factor $2$ in $3a_{k}+1$, where $a$ is the initial number given, and $k$ runs in $\mathbb{N}$.
With this we can redefine the sequence with
\[ a_{k + 1} = \frac{3a_{k}+1}{2^{\val(a;k+1)}} = \rel_{2}(3a_{k}+1), \]
which ensures the $a_k$ always be odd.
With the definition $\mathcal{A}_{k} = \sum_{j=0}^{k-1}\val(a;j+1)$, we have
\[ a_k \cdot 2^{\mathcal{A}_k} = a \cdot 3 ^ k + \sum_{i = 0}^{k - 1}2^{\mathcal{A}_i}\cdot3^{k-i-1}. \label{equation:collatz-inductioon-3-2-1}\]
It means $a \equiv -3^{-1} \cdot \sum_{i = 0}^{k - 1}2^{\mathcal{A}_i}\cdot3^{-i} \pmod{2^{\mathcal{A}_k}}$.

Notice that $a$ is a positive finite number but $2^{\mathcal{A}_k}$ increase to infinite, that means $\sum_{i = 0}^{k - 1}2^{\mathcal{A}_i}\cdot3^{-i}$ should be a finite odd times of $3^{-1}$ under the modulo, which convergences to $1/3$ in $\mathbb{Z}_2$.
And this congruence equation system has one obvious solution: $a = 1$, that $\mathcal{A}_k = 2k$, holds
\[ -3^{-1} \cdot \sum_{i = 0}^{k - 1}2^{\mathcal{A}_i}\cdot3^{-i} = - 3^{-k}\cdot 4^k + 1 \equiv 1 \pmod{2^{2k} = 4^{k}}, \]
follows the Collatz conjecture.
We can see that each number generates a sequence $\mathcal{A}_{k}$ such that the above congruence equation holds.
So it is the time to explore what the arguments of these equations has a stability solution.

At first we define an auxity notation:
\begin{definition}~\label{def:exclude-beta-times}
    For the number field $\mathbb{K}$ ($\mathbb{N}$, $\mathbb{Q}$, or $p$-adic integer $\mathbb{Z}_{p}$), we denote that $\widehat{\mathbb{K}}_{{\beta}} = \{a \in\mathbb{K} \mid \beta \nmid a\}$.
\end{definition}
Here the divisibility follows \cref{def:extend-divisibility}, and any divisibility appearing in this text follows this definition.
With these we have
\begin{definition}[Collatz Sequence and Collatz Reduce Sequence]~\label{definition:collatz-sequence}
    Given arguments $\alpha, \beta\in \mathbb{N}_{\ge2}$ and $\gamma \in\mathbb{N}_{\ge 1}$ that $\gcd(\beta, \alpha\gamma) = 1$.\newline
    For number $a \in \widehat{\mathbb{N}}_{{\beta}}$, define the sequence $\left\{a_{k}\right\}_{k = 0}^{\infty} \subset \widehat{\mathbb{N}}_{{\beta}}$ named Collatz sequence and $\{\val(a;k)\}_{k=0}^{\infty}\subset\mathbb{N}$ named Collatz reduce sequence with
    \begin{align*}
        &a_0=a, \\
        &\val(a;k) = 0, \\
        &a_{k+1} = \frac{\alpha\cdot a_{k} + \gamma}{\beta^{\val(a;k + 1)}} = \rel_{\beta}(\alpha\cdot a_{k} + \gamma), \\
        &\val(a;k+1) = \val_{\beta}(\alpha\cdot a_k + \gamma).
    \end{align*}
\end{definition}
Actually, this definition is exactly follow the Syracuse function\cite{wikipediaCollatzConjecture}. 
\begin{definition}~\label{def:collatz-accumulation}
    Given a sequence of number $\{v_{i}\}\subset\mathbb{N}$, finite or infinite terms. $\forall m, n \in \mathbb{N}$ that $n \ge m$, denote $\mathcal{V}_{n}^{m} = \sum_{i=m}^{n-1}v_{i}$ named accumulation from $m$ to $n$.
    Specially when $m = 0$, we omit the superscript, that $\mathcal{V}_n = \mathcal{V}_n^{0}$.
\end{definition}
Normally, we will always lowercase latter like $u$, $v$, $w$, $\cdots$ to represent elements of some sequences or vectors.
\begin{proposition}
    $\forall n \ge k \ge m$ there is $\mathcal{V}_n^m = \mathcal{V}_n^{k} + \mathcal{V}_k^{m}$.
\end{proposition}

Given a number $a \in \widehat{\mathbb{N}}_{{\beta}}$, we can derive a Collatz sequence $\left\{a_{k}\right\}_{k = 0}^{\infty}$ and Collatz reduce sequence $\{\val(a;k)\}_{k=0}^{\infty}$. For the Collatz reduce sequence, we can define its accumulation, and we call it Collatz accumulation $\mathcal{A}$.
And to avoid repeating the explanation each time, unless otherwise specified, we will use the uppercase calligraphic font to represent the Collatz accumulation of the corresponding letters by default.
\begin{proposition}[Uniqueness]~\label{proposition:uniqueness}
If $a = b$, then $\forall k \in \mathbb{N}$, there is $a_k = b_k$ and $\val(a;k) = \val(b;k)$ and $\mathcal{A}_k = \mathcal{B}_k$.
\end{proposition}
% \begin{proof}
%     First, we have
%     \begin{align*}
%         a = b &\Rightarrow a_0 = b_0, \\
%         a_1 \cdot \beta^{\val(a;1)} &= b_1 \cdot \beta^{\val(b;1)}, \\
%         \mathcal{A}_1 = \val(a;1) &= \val(b;1) = \mathcal{B}_1.
%     \end{align*}
%     Second, if $a_k = b_k$ and $\mathcal{A}_k = \mathcal{B}_k$, we have
%     \[a_{k+1} \cdot \beta^{\val(a;k + 1)} = \alpha\cdot a_{_k} + \gamma = \alpha\cdot b_{_k} + \gamma = b_{k+1} \cdot \beta^{\val(b;k + 1)}.\]
%     That means
%     \[a_{k+1} = b_{k+1} \cdot \beta^{\val(b;k + 1) - \val(a;k + 1)}. \]
%     But there is $a_{k+1} \in \widehat{\mathbb{N}}_{{\beta}}$ as definition that $\beta \nmid a_{k+1}$, which means
%     \[\val(a;k + 1) = \val(b;k + 1), \]
%     and then
%     \begin{align*}
%         a_{k+1} &= b_{k+1}, \\
%         \mathcal{A}_{k+1} = \mathcal{A}_{k} + \val(a;k+1) &= \mathcal{B}_{k} + \val(b;k+1) = \mathcal{B}_{k+1}.
%     \end{align*}
%     Sum up, we proved this proposition.
% \end{proof}
\begin{proposition}
    If $\exists t > 0$ that $a_{t} = a_0 = a$, 
    then $\forall k > 0$, and $s = 1, 2, 3, \cdots, t$, there are
    $a_{kt+s} = a_s$ and $\mathcal{A}_{kt + s}^{kt} = \mathcal{A}_{s}$.
\end{proposition}
\begin{proposition}~\label{proposition:trivial-periodic-sequence}
    If there is $\Delta \in \mathbb{N}_{\ge1}$ holds $(\beta^{\Delta} - \alpha) \mid \gamma$, then let $a_0 = \frac{\gamma}{\beta^{\Delta} - \alpha}$, we have $\forall k \in \mathbb{N}$, $a_k = a_0$ and $A_k = k\cdot \Delta$.
\end{proposition}
\begin{proposition}~\label{proposition:equivalent-between-collatz-conjecture-and-our-definition}
Let arguments $(\alpha, \beta, \gamma) = (3, 2, 1)$.
For all odd number $a$, let $a'_{i_k}$ be the $k$-th odd number in \cref{conjecture:original-conjecture} and $a_k$ be the $k$-th element of Collatz sequence of $a$, then $\forall k \in \mathbb{N}$,  we have $a_{k} = a'_{i_{k}}$.
\end{proposition}
% \begin{proof} Trivial. \end{proof}
% \begin{proof}
%     Let $a'_{i}$ be the $i$-th number in \cref{conjecture:original-conjecture}, and $a'_{i_k}$ be the $k$-th odd number in it.\newline
%     First, we have $a_{0} = a = a'_0 = a'_{i_0}$;\newline
%     Second, if $a_{k} = a'_{i_k}$. Suppose $s_{k+1} = \val_{\beta}(3a'_{i_k} + 1)$. We have
%     \[ s_{k+1} = \val_{\beta}(3a_k + 1) = \val(a;k+1). \]
%     For $a'_{i_k}$, we have
%     \begin{align*}
%         a'_{i_k+1} &= (3a'_{i_k} + 1) / 2 \in 2\mathbb{N},   \\
%         a'_{i_k+2} &= (3a'_{i_k} + 1) / 2^2 \in 2\mathbb{N}, \\
%         &\cdots \\
%         a'_{i_k+s_k} &= (3a'_{i_k} + 1) / 2^{s_{k+1}}  \in \widehat{\mathbb{N}}_{2}.
%     \end{align*}
%     That means $a'_{i_k+s_k}$ is the $(k+1)$-th odd number,  in other words, $a'_{i_k+s_k} = a'_{i_{k+1}}$.
%     Hence we have
%     \begin{align*}
%         a_{k+1} &= (3a_k + 1) / 2^{\val(a;k+1)} \\
%         &= (3a'_{i_k} + 1) / 2^{s_{k+1}} \\
%         & = a'_{i_k+s_k} \\
%         &= a'_{i_{k+1}}.
%     \end{align*}
%     Sum up, we proved this proposition.
% \end{proof}
\cref{proposition:equivalent-between-collatz-conjecture-and-our-definition} ensures our definition of Collatz sequence is equivalent with the iterative function in \cref{conjecture:original-conjecture}.
Because for even numbers, we can easily figure out that every even number will reduce to an odd number at first.
In that situation, the proposition holds as well.
With this proposition, we can describe the Collatz Conjecture with the Collatz sequence:

\begin{conjecture}[Rewritten Collatz Conjecture]~\label{definition:rewritten-conjecture}
With arguments $(\alpha, \beta, \gamma) = (3, 2, 1)$,
for any odd number $a \in \widehat{\mathbb{N}}_{2}$, it's Collatz sequence will finally reach to $1$.
\end{conjecture}

This article involves many terms like $\textstyle \sum_{i = 0}^{k - 1}\beta^{\mathcal{A}_i}\cdot \alpha^{k-i-1}$. We first define a notation to represent them:
\begin{definition}~\label{definition:collatz-complement}
    Given a sequence of number $\{v_{i}\}\subset\mathbb{N}$, finite or infinite terms, we have $\mathcal{A}_{n}^{m}$. Now we define $\mathfrak{A}_n^{m}$ called Collatz complement of $\{v_{i}\}$ at $n$ relative to $m$ with
    \[ \mathfrak{A}_n^{m} = \sum_{i = m}^{n - 1}\beta^{\mathcal{A}_i^{m}}\cdot\alpha^{n-i-1} \in \mathbb{N}. \]
    And define $\mathfrak{a}_n^{m}$ called reduced Collatz complement of $a$ at $n$ relative to $m$ with
    \[ \mathfrak{a}_n^{m} = \frac{\mathfrak{A}_n^{m}}{\alpha^{n-1}} = \sum_{i = m}^{n - 1}\frac{\beta^{\mathcal{A}_i^{m}}}{\alpha^{i}} \in \mathbb{Q}. \]
    Also, they can be omitted the superscript $m$ when $m = 0$ that $\mathfrak{A}_n = \mathfrak{A}_n^{0}$, and $\mathfrak{a}_n = \mathfrak{a}_n^{0}$.
\end{definition}

If $\{v_{i}\}$ is the Collatz reduces of a number, to avoid repeating the explanation each time, without further explanation, we will use the same letter of the Collatz accumulation but uppercase fraktur font to represent Collatz complement, and the same letter but lowercase fraktur font to represent reduced Collatz complement.
In that case, we also say that $\mathfrak{A}$ and $\mathfrak{a}$ are of number $a$.

For the Collatz complement, we have two basic propositions:
\begin{proposition}~\label{proposition:collatz-complement-basic-prop1}
$\mathfrak{A}_{n}^{n} = 0$, and $\mathfrak{A}_{n+1}^{n} = 1$.
And $\mathfrak{a}_{n}^{n} = 0$, and $\mathfrak{a}_{n+1}^{n} = \frac{1}{\alpha^{n}}$.
\end{proposition}
% \begin{proposition}
%     $\mathfrak{A}_{n+1} = \alpha \cdot \mathfrak{A}_{n} + \beta^{\mathcal{A}_n}$. And $\mathfrak{a}_{n+1} = \mathfrak{a}_{n} + \frac{\beta^{\mathcal{A}_n}}{\alpha^{n}}$.
% \end{proposition}
\begin{proposition}~\label{proposition:collatz-complement-basic-prop2}
    $\forall n \ge k \ge m$, we have
    \begin{align*}
        \mathfrak{A}_n^{m} &= \beta^{\mathcal{A}_k^{m}}\cdot\mathfrak{A}_n^{k} + \mathfrak{A}_k^m\cdot\alpha^{n-k}, \\
        \mathfrak{a}_n^{m} &= \beta^{\mathcal{A}_k^{m}}\cdot\mathfrak{a}_n^{k} + \mathfrak{a}_k^m. 
    \end{align*}
\end{proposition}
% \begin{proof}
% \begin{align*}
%     \mathfrak{A}_n^{m} =&\sum_{i = m}^{n - 1}\beta^{\mathcal{A}_i^{m}}\cdot\alpha^{n-i-1} \\
%     =&\ \sum_{i = k}^{n - 1}\beta^{\mathcal{A}_i^{m}}\cdot\alpha^{n-i-1} + \sum_{i = m}^{k - 1}\beta^{\mathcal{A}_i^{m}}\cdot\alpha^{n-i-1} \\
%     =&\ \sum_{i = k}^{n - 1}\beta^{\mathcal{A}_k^{m}}\cdot\beta^{\mathcal{A}_i^{k}}\cdot\alpha^{n-i-1} + \sum_{i = m}^{k - 1}\beta^{\mathcal{A}_i^{m}}\cdot\alpha^{k-i-1}\cdot\alpha^{n-k} \\
%     =&\ \beta^{\mathcal{A}_k^{m}}\cdot\sum_{i = k}^{n - 1}\beta^{\mathcal{A}_i^{k}}\cdot\alpha^{n-i-1} + \sum_{i = m}^{k - 1}\beta^{\mathcal{A}_i^{m}}\cdot\alpha^{k-i-1}\cdot\alpha^{n-k} \\
%     =&\ \beta^{\mathcal{A}_k^{m}}\cdot\mathfrak{A}_n^{k} + \mathfrak{A}_k^m\cdot\alpha^{n-k}, \\
%     \mathfrak{a}_n^{m} =&\sum_{i = m}^{n - 1}\frac{\beta^{\mathcal{A}_i^{m}}}{\alpha^{i}} \\
%     =&\ \sum_{i = k}^{n - 1}\frac{\beta^{\mathcal{A}_i^{m}}}{\alpha^{i}} + \sum_{i = m}^{k - 1}\frac{\beta^{\mathcal{A}_i^{m}}}{\alpha^{i}} \\
%     =&\ \beta^{\mathcal{A}_k^{m}}\cdot\sum_{i = k}^{n - 1}\frac{\beta^{\mathcal{A}_i^{k}}}{\alpha^{i}} + \sum_{i = m}^{k - 1}\frac{\beta^{\mathcal{A}_i^{m}}}{\alpha^{i}} \\
%     =&\ \beta^{\mathcal{A}_k^{m}}\cdot\mathfrak{a}_n^k + \mathfrak{a}_k^m.
% \end{align*}
% \end{proof}

Actually, we can formally define the Collatz complement only from \cref{proposition:collatz-complement-basic-prop1} and \cref{proposition:collatz-complement-basic-prop2}.
More essentially, the two propositions are exactly the basic rules of a quasi $(\sigma, \tau)$-derivation of \cref{sec:a-kind-of-derivation-on-groupoid}, which in other words, the Collatz complement is a kind of $(\sigma, \tau)$-derivation on groupoid (\cref{def:quasi-derivation-on-groupoid}).

According to \cref{proposition:uniqueness} and \cref{cor:co-uniqueness}, we have
\begin{proposition}[Uniqueness]~\label{proposition:uniquess-of-Collatz-complement}
    Given $\mathfrak{V}$ for $\{v_i\}$ and $\mathfrak{U}$ for $\{u_i\}$. If $\forall i\ge 0$, there is $v_i = u_i$, then $\forall n \ge m \ge 0$, there are $\mathfrak{V}_n^{m} = \mathfrak{U}_n^{m}$ and $\mathfrak{v}_n^{m} = \mathfrak{u}_n^{m}$.
\end{proposition}
\begin{proposition}~\label{prop:complement-equiv-1-when-modulo}
    Suppose $m\ne n$, we have $\mathfrak{V}_n^{m} \equiv 1 \pmod{\alpha}$. And if $v_{1} \ne 0$, then we also have $\mathfrak{V}_n^{m} \equiv 1 \pmod{\beta}$.
\end{proposition}
Now we are ready to explore the properties of Collatz sequences.
But before that, I need to point out that there are many symbols in number theory.
These symbols, while simple, differ slightly from those in some other papers.
These symbols are defined in \cref{sec:number-theory-tools}. Before proceeding, please ensure you are familiar with these symbols first.
